% abstract.tex -- standalone abstract, \input into paper.tex and also kept
% separate for submission portals that require the abstract as its own file.
%
% Rewritten against the LOCKED Introduction/Methods/Results/Discussion
% (manuscript_blueprint.md status: all four sections locked). Replaces the
% pre-lock draft written before E2.6/E2.7 were finalized in prose. Ends on
% the plain slogan wording (matches Introduction's instance), per the
% decision in manuscript_blueprint.md's slogan amendment: Discussion's
% "The same information..." callback stays unique to the close of the paper.

Distance-based reliability estimation assumes that a representation's
geometry reflects its trustworthiness, yet this assumption is rarely tested
under training interventions that reshape representation geometry directly.
We audit this assumption under domain-adversarial representation learning, using a disentanglement dose-response ladder. Three
checkpoint families share the same architecture and a 16-dimensional
representation, differing only in orthogonality strength
($\lambda_{orth} = 0, 1, 5$). Representation geometry changed
substantially with disentanglement strength: condition number shifted by
two orders of magnitude (Kendall's $\tau=0.84$, exact $p=2.8\times10^{-5}$).
This geometric change was not accompanied by any improvement in reliability
estimation: Mahalanobis-distance AUROC (ISIC-test vs.\ PAD-UFES) remained
flat and below chance ($\approx 0.40$) at every disentanglement level, with
no significant association to any of five geometry metrics tested. The same
failure was shared by cosine-to-centroid and pooled $k$-nearest-neighbor
scorers, as well as three additional, non-distance-based scorers (an
energy-based confidence score, Virtual-Logit Matching, and a kernel density
estimator) -- eight scorers spanning distinct scoring principles, seven of
which converged on the same result; one (the energy-based score) showed an
isolated upward trend with disentanglement strength that we report but do
not treat as evidence against the overall pattern. A supervised probe with no access to the
training objective nonetheless recovered domain membership from the
identical embeddings at 0.72--0.81 AUROC across every disentanglement level,
showing the relevant information was not absent from the representation.
Collectively, these findings indicate that evaluating a learned representation by
classification performance alone can overlook whether the information it
contains is organized in a form downstream reliability estimators can
actually use. Information can remain decodable while becoming largely inaccessible
to non-probing reliability estimators.
